\documentclass[12pt]{article}
\usepackage[utf8]{inputenc}
\usepackage[T1]{fontenc}
\usepackage{hyperref}
\title{W\`e Western Language Resources for K-12 Curriculum Development}
\author{Compiled by Claude Sonnet 5 (claude-sonnet-5) for Andr\'as Kornai}
\date{}
\begin{document}
\maketitle
\section{Introduction}
W\`e Western (ISO 639-3 wec) belongs to the Kru group in C\^ote d'Ivoire. The W\`e (Gu\'er\'e/Wob\'e) cluster has dialect groups that have been described together. My searches found little dedicated to wec itself. I found no dictionary, primer or textbook specific to it, and the resources below are general to the W\`e/Gu\'er\'e cluster.
\section{Primary Educational Resources}
None found in search results.
\section{Linguistic and Reference Materials}
\begin{itemize}
\item Overview of the Guere language, with a phonology sketch and classification: \url{https://en.wikipedia.org/wiki/Guere_language}
\item Comparative study on the unity and diversity of the W language cluster of C\^ote d'Ivoire (Journal of West African Languages): \url{https://journalofwestafricanlanguages.org/downloads/send/93-volume3302/458-unite-et-diversite-du-w-de-cote-d-ivoire}
\item Omniglot page on Gu\'er\'e (speaker and literacy estimates): \url{https://www.omniglot.com:443/writing/guere.htm}
\end{itemize}
A search snippet cited Paradis (1983), \emph{Description phonologique du gu\'er\'e}. I did not see a URL for it.
\section{Religious and Community Resources}
None found in search results.
\section{Digital and Audio Resources}
None found in search results.
\section{Curriculum Development Context}
Direct evidence of school use is lacking. A curriculum would have to start from basic documentation, with community and linguist involvement. Care is needed because Gu\'er\'e and Wob\'e varieties differ.
\end{document}